% =============================================================================
% Guía interna del grupo — Tarea 1: funcionamiento del código y guion de la
% presentación (presentacion_tarea1.pptx). No es parte de la entrega.
% Compilar con pdfLaTeX; no necesita imágenes.
% =============================================================================

\documentclass[11pt]{article}
\usepackage[T1]{fontenc}
\usepackage[utf8]{inputenc}
\usepackage[spanish, es-tabla, es-nodecimaldot, es-noshorthands]{babel}
\usepackage{amsmath,amssymb}
\usepackage{xcolor}
\usepackage{enumitem}
\usepackage{booktabs}
\usepackage{array}
\usepackage[margin=2.5cm]{geometry}
\usepackage{listings}
\usepackage{hyperref}
\hypersetup{colorlinks=true, linkcolor=black, urlcolor=blue}
\definecolor{codegreen}{rgb}{0,0.6,0}
\definecolor{codegray}{rgb}{0.5,0.5,0.5}
\definecolor{backcolour}{rgb}{0.96,0.96,0.96}
\definecolor{guion}{rgb}{0.93,0.95,1.0}

\lstset{
    backgroundcolor=\color{backcolour},
    commentstyle=\color{codegreen},
    keywordstyle=\color{blue},
    numberstyle=\tiny\color{codegray},
    basicstyle=\ttfamily\footnotesize,
    breaklines=true,
    keepspaces=true,
    numbers=left,
    numbersep=5pt,
    showstringspaces=false,
    tabsize=2,
    frame=single,
    rulecolor=\color{black!10},
    literate={á}{{\'a}}1 {é}{{\'e}}1 {í}{{\'i}}1 {ó}{{\'o}}1 {ú}{{\'u}}1 {ñ}{{\~n}}1
}

\setlist{itemsep=2pt,topsep=3pt}

\newcommand{\code}[1]{\texttt{#1}}

% Bloque de "lo que se dice" en el guion.
\newcommand{\decir}[1]{%
  \par\noindent\colorbox{guion}{\parbox{\dimexpr\linewidth-2\fboxsep}{\small\textit{#1}}}\par\medskip}

\title{Tarea 1 — Guía del código y guion de la presentación\\[4pt]\large Uso interno del grupo}
\author{}
\date{1 de octubre de 2026}

\begin{document}
\maketitle
\tableofcontents
\bigskip


% =============================================================================
\part{Funcionamiento del código}
% =============================================================================

\section{Qué hace la tarea, en una frase}

Se toma una traza real de tráfico de Internet (MAWI, 15 minutos, unos 123 millones de paquetes), se le inyectan dos ataques sintéticos (un DDoS y un scan) y se intenta detectarlos estimando, para cada ventana de los últimos $W=60$\,s que avanza cada $p=10$\,s, cuántos paquetes tiene la IP sospechosa. La estimación se hace con dos estructuras probabilísticas de memoria pequeña, \textbf{Count-Min Sketch (CMS)} y \textbf{CountSketch (CS)}, y se compara contra un conteo exacto.

Un paquete es \emph{heavy hitter} en la ventana $j$ si su frecuencia supera $\lceil\varphi N_j\rceil$, con $\varphi=0{,}01$ (el 1\,\% de los paquetes de la ventana). Detectar el ataque es que la IP de la víctima (DDoS) o del atacante (scan) cruce ese umbral.

\textbf{No hay interfaz gráfica.} Todos los programas son de consola: leen archivos binarios, imprimen un resumen en la terminal y escriben CSV. El resultado visible son las figuras PNG y las tablas Markdown de \code{codigo\_entregado/resultados/}.

\section{Flujo de datos entre programas}

\begin{enumerate}
    \item \textbf{\code{201812031400.pcap.gz}} (captura MAWI original, $\sim$3\,GB) \\
          $\downarrow$ \code{pcap2bin} (C++, del curso)
    \item \textbf{\code{traza.bin}}: 123\,383\,971 registros de 24\,bytes, sin ataques \\
          $\downarrow$ \code{inject\_attack.py} (Python, del curso), semilla 42
    \item \textbf{\code{traza\_ddos.bin}}, \textbf{\code{traza\_scan.bin}} (y sus versiones \code{\_bajo}) + \code{gt\_*.json} con el \emph{ground truth} \\
          $\downarrow$ en paralelo:
    \begin{itemize}
        \item \code{exact\_hh} (C++, del curso) $\rightarrow$ \code{exact\_<exp>.csv}: frecuencia exacta por ventana.
        \item \code{sliding\_sketch} (C++, \textbf{nuestro}) $\rightarrow$ \code{sk\_<exp>\_s<semilla>.csv}: frecuencia estimada por CMS y CS.
    \end{itemize}
    \item $\downarrow$ \code{analisis.py} (Python, nuestro) une ambos CSV
    \item \textbf{\code{resultados/}}: figuras \code{fig\_*.png} y tablas \code{resumen.md}, \code{validacion.md}, \code{delta\_extremos.md}.
\end{enumerate}

Todo lo anterior lo orquesta \code{run\_experimentos.sh}, que llama a \code{preparar\_trazas.sh} para los pasos 1–3.

\subsection*{Formato de un registro (24 bytes)}

\begin{center}
\begin{tabular}{lll}
\toprule
Campo & Tipo & Significado \\
\midrule
\code{ts\_us} & \code{uint64} & instante en microsegundos desde 1970 \\
\code{src}, \code{dst} & \code{uint32} & IP de origen y de destino \\
\code{sport}, \code{dport} & \code{uint16} & puertos \\
\code{len} & \code{uint16} & bytes del paquete \\
\code{proto}, \code{flags} & \code{uint8} & protocolo y banderas (el bit \code{0x10} marca paquetes sintéticos) \\
\bottomrule
\end{tabular}
\end{center}

Leer un arreglo plano de registros de tamaño fijo es mucho más rápido que volver a parsear el pcap en cada experimento.

\section{Cada programa}

\subsection{\code{pcap2bin.cpp} (entregado por el curso)}
Lee el pcap por la entrada estándar y escribe un registro de 24\,B por cada paquete IPv4. Sólo se modificó para que compile fuera de Windows (\code{<io.h>} y \code{\_setmode} quedaron bajo \code{\#ifdef \_WIN32}); la conversión es idéntica.

\subsection{\code{inject\_attack.py} (entregado por el curso)}
Genera los paquetes del ataque y los mezcla, ordenados por tiempo, con la traza base. Escribe el JSON con los parámetros reales del ataque.
\begin{itemize}
    \item \textbf{DDoS}: 10\,000 paquetes/s durante 30\,s desde 4\,000 IPs distintas hacia la víctima \code{163.210.30.13}. Se suman 300\,000 paquetes.
    \item \textbf{Scan}: 8\,000 paquetes/s durante 30\,s desde \code{198.18.0.7} hacia 60\,000 destinos. Se suman 240\,000 paquetes.
    \item Variantes \code{\_bajo}: los mismos ataques a 3\,000 paquetes/s, para que el error de los sketches se note.
\end{itemize}
Ambos comienzan en $t=300$\,s. La víctima del DDoS se fija explícitamente porque, en modo automático, el script la elige con un ordenamiento no estable y distintas versiones de numpy elegían IPs distintas.

\subsection{\code{exact\_hh.cpp} (entregado por el curso)}
La referencia exacta: con una tabla hash cuenta todas las claves de cada ventana. Con \code{--query IP --out-query} entrega, por ventana, $N_j$, el umbral, la frecuencia exacta $f_j(x)$, si es heavy hitter y el cambio exacto $\Delta f_j$. Usa \code{mmap}, por eso en Windows se compila y corre en WSL.

\subsection{\code{elegir\_claves.py} (nuestro)}
Para validar los sketches \emph{sin} ataque, elige en la primera ventana de \code{traza.bin} cuatro grupos de 10 IPs: las 10 más frecuentes, y las 10 más cercanas a frecuencia 1000, 100 y 10. Lo hace por origen (\code{claves\_src.txt}) y por destino (\code{claves\_dst.txt}).

\subsection{\code{sliding\_sketch.cpp} (nuestro: la implementación de la tarea)}
Recorre la traza una sola vez y mantiene, para cada ancho $w$ pedido, un CMS y un CS sobre la ventana deslizante. En cada evaluación escribe una fila CSV por IP consultada y por ancho. Se explica en detalle en la sección~\ref{sec:sliding}.

\subsection{\code{analisis.py} (nuestro)}
Une \code{sk\_*.csv} con \code{exact\_*.csv} por ventana e IP, y calcula:
\begin{itemize}
    \item \textbf{MRE} (error relativo medio) sobre las 8 ventanas afectadas por el ataque.
    \item \textbf{Latencia}: segundos desde el inicio del ataque hasta la primera ventana en que se declara heavy hitter, para el exacto y para cada sketch.
    \item \textbf{Falsos positivos y negativos} contra la decisión de \code{exact\_hh}.
    \item \textbf{Memoria} de contadores: $7\cdot d\cdot w\cdot 4$\,B.
    \item Error de $\Delta f$ y sus extremos; error sin ataque por grupo de claves.
\end{itemize}
Produce \code{fig\_<ataque>.png} (frecuencia), \code{fig\_delta\_<ataque>.png} (cambio) y las tablas.

\subsection{Scripts y \code{Makefile}}
\begin{itemize}
    \item \code{preparar\_trazas.sh}: verifica requisitos, compila, descarga el pcap (con reintentos), lo convierte, inyecta los ataques y comprueba el número exacto de registros de cada traza.
    \item \code{run\_experimentos.sh}: el único script que hay que correr. Es idempotente (omite lo que ya existe; \code{FORCE=1} rehace todo).
    \item \code{Makefile}: compila los tres programas C++.
\end{itemize}

\section{\code{sliding\_sketch.cpp} en detalle}\label{sec:sliding}

\subsection{Funciones hash}
Se usa la familia de Carter–Wegman $h(x)=(a x+b)\bmod(2^{61}-1)$, con $a,b$ sacados de un generador \code{mt19937\_64} con la semilla de hash (42, 7 o 1234). Cada una de las $d=5$ filas tiene dos funciones independientes:
\begin{itemize}
    \item \textbf{posición}: $h_i(x)\bmod w$, la columna del contador;
    \item \textbf{signo} (sólo CS): $s_i(x)=\pm1$ según el bit menos significativo de otra función de la familia.
\end{itemize}
Las mismas funciones se usan para todos los anchos; sólo cambia el módulo.

\subsection{Los dos sketches}
Ambos son una matriz de $d\times w$ contadores. Se diferencian en sólo dos cosas, encapsuladas en \code{struct CountMin} y \code{struct CountSketch}:

\begin{center}
\begin{tabular}{p{3.2cm}p{5.4cm}p{5.4cm}}
\toprule
 & \textbf{CMS} & \textbf{CS} \\
\midrule
Actualizar con $x$ & $A[i][h_i(x)] \mathrel{+}= 1$ & $A[i][h_i(x)] \mathrel{+}= s_i(x)$ \\
Estimar $f(x)$ & $\min_i A[i][h_i(x)]$ & $\operatorname{mediana}_i\, s_i(x)\,A[i][h_i(x)]$ \\
Tipo de contador & \code{uint32} & \code{int32} \\
Error & sólo sobreestima (colisiones suman) & sin sesgo, error en ambos sentidos \\
\bottomrule
\end{tabular}
\end{center}

\subsection{La ventana deslizante: \code{SlidingWindow<Sketch>}}

Es una plantilla que no sabe qué sketch contiene; sólo le pide el signo de la actualización y el estimador. Guarda:
\begin{itemize}
    \item un \textbf{anillo} de $m=W/p=6$ sub-sketches $S$, uno por subventana de 10\,s;
    \item el \textbf{agregado} $A=S_1+\dots+S_6$, que representa la ventana completa de 60\,s;
    \item \code{dA}, el sketch del cambio entre la ventana anterior y la actual.
\end{itemize}

La subventana $q$ cubre $(t_0+(q-1)p,\ t_0+qp]$ y vive en la ranura $(q-1)\bmod 6$. Los tiempos se tratan como enteros en microsegundos para que no haya errores de redondeo en los bordes.

\begin{lstlisting}[language=C++,caption={Las tres operaciones de la ventana (simplificado)}]
// Cada paquete: se suma en su sub-sketch Y en el agregado.
void add(int slot, col, sg) {
    for (int i = 0; i < d; ++i) {
        k = i * w + col[i];
        ring[slot][k] += sign(sg[i]);
        A[k]          += sign(sg[i]);
    }
}
// Al abrir una subventana nueva en una ranura ocupada: se resta lo viejo.
void expire(int slot) {
    for (k = 0; k < d * w; ++k) {
        dA[k] = -ring[slot][k];
        A[k] -= ring[slot][k];
        ring[slot][k] = 0;            // la ranura queda lista para reutilizarse
    }
}
// Al cerrar la subventana nueva: dA = S_nueva - S_vieja.
void close_delta(int slot) {
    for (k = 0; k < d * w; ++k) dA[k] += ring[slot][k];
}
\end{lstlisting}

\paragraph{Uso de la linealidad.} Un sketch es una función lineal del flujo: el sketch de la unión de dos conjuntos de paquetes es la suma de sus sketches, y el de una diferencia es la resta. Por eso
\[
A_j = A_{j-1} - S_{\text{sale}} + S_{\text{entra}}, \qquad \Delta A_j = A_j - A_{j-1} = S_{\text{entra}} - S_{\text{sale}}.
\]
Cada rotación toca $2dw$ contadores, \textbf{sin importar cuántos paquetes tenga la ventana} (más de 8 millones). Sin linealidad habría que reconstruir el sketch de 60\,s reprocesando todos sus paquetes cada 10\,s. Esto funciona igual para CMS y CS; la única diferencia es que en CMS los contadores de $\Delta A$ pueden ser negativos, y entonces el mínimo deja de tener sentido (ver \S\ref{sec:delta}).

\subsection{El ciclo principal}
\begin{enumerate}
    \item Lee la traza en bloques de 65\,536 registros.
    \item Para cada paquete calcula su subventana $q=\lceil(t-t_0)/p\rceil$. Si $q$ avanzó, cierra la subventana anterior (y evalúa si se completó una ventana) y abre la siguiente, expirando la que ocupaba su ranura. Esto también recorre subventanas vacías.
    \item Calcula las $d$ posiciones y signos de la clave (\code{src} para el scan, \code{dst} para el DDoS) y la agrega a los CMS y CS de todos los anchos a la vez.
    \item Mantiene $N_j$ exacto con un anillo de 6 contadores escalares; con él se calcula el umbral $\lceil\varphi N_j\rceil$.
\end{enumerate}

La primera evaluación ocurre en $t_0+60$\,s, cuando las seis ranuras ya están llenas; después, una cada 10\,s, hasta 84 ventanas.

\paragraph{Autoverificación.} Con \code{--verify exact\_<exp>.csv}, el programa compara su $N_j$ con el de \code{exact\_hh} en cada ventana y termina con código 3 si alguna difiere. Así se garantiza que ambos programas alinean las ventanas igual; en todas las corridas coincidieron las 84 ventanas.

\subsection{Salida}
Por la terminal (stderr) imprime un resumen: traza, parámetros, ventanas evaluadas, tiempo de la pasada ($\sim$14\,s) y memoria por ancho. El CSV tiene una fila por ventana, IP y ancho:

\begin{center}
\small
\begin{tabular}{ll}
\toprule
Columna & Significado \\
\midrule
\code{win}, \code{t\_rel\_s} & número de ventana y segundo en que termina \\
\code{N}, \code{threshold} & paquetes en la ventana y umbral $\lceil\varphi N\rceil$ \\
\code{cms\_f}, \code{cms\_hh} & estimación CMS y si supera el umbral \\
\code{cs\_f}, \code{cs\_hh} & estimación CS (truncada a 0) y si supera el umbral \\
\code{cms\_med\_delta}, \code{cs\_delta} & cambio $\Delta f$ estimado (CMS-mediana y CS) \\
\bottomrule
\end{tabular}
\end{center}

\subsection{El cambio $\Delta f$ y la CMS-mediana}\label{sec:delta}
$\Delta f_j(x)=f_j(x)-f_{j-1}(x)$ indica cuándo empieza el ataque (salto positivo grande) y cuándo su tráfico sale de la ventana (salto negativo). CS lo estima con su estimador normal sobre \code{dA}, porque ya trabaja con valores con signo. Para CMS se usa la mediana de los contadores (sin signo de hash), llamada \textbf{CMS-mediana}: es una variante experimental que pierde la garantía del CMS, que depende de que los contadores sean no negativos.

\section{Cómo correrlo y qué se ve}

\begin{lstlisting}[language=bash,numbers=none]
cd codigo_entregado
bash run_experimentos.sh            # FORCE=1 bash run_experimentos.sh rehace todo
\end{lstlisting}

Requiere \code{g++} con C++17, Python 3 con \code{numpy}, \code{pandas} y \code{matplotlib}, y en Windows WSL para \code{exact\_hh}. Si las trazas no existen, primero descarga el pcap ($\sim$3\,GB). Con las trazas listas tarda unos 10 minutos. En la terminal se ven los pasos y los resúmenes de cada programa; al final, en \code{resultados/}:
\begin{itemize}
    \item \code{fig\_ddos.png}, \code{fig\_scan.png}: frecuencia exacta contra estimada entre 240 y 390\,s, con el umbral.
    \item \code{fig\_delta\_ddos.png}, \code{fig\_delta\_scan.png}: $\Delta f$ exacto contra estimado.
    \item \code{resumen.md}: MRE, memoria, latencia, falsos positivos y negativos.
    \item \code{validacion.md}: errores sin ataque por grupo de claves.
\end{itemize}

% =============================================================================
% =============================================================================
\part{Guion de la presentación}
% =============================================================================

\section{Plan general}

El guion sigue, en el mismo orden, las 9 diapositivas de \code{codigo\_entregado/presentacion\_tarea1.pptx}. Cubre lo que pide el enunciado: \textbf{diseño de la ventana deslizante, uso de la linealidad, resultados para ambos ataques y conclusiones de la comparación CMS frente a CS}. Las cifras son las mismas que muestran las diapositivas y sus notas del presentador.

El texto sombreado es lo que se propone decir. Son unas 790 palabras: entre 6 y 7 minutos a ritmo normal de exposición (110 a 130 palabras por minuto). Los tiempos de la tabla suman \textbf{8 minutos} porque incluyen pausas para señalar los gráficos; el resto, hasta 10, queda de margen para cambios de diapositiva y traspasos. Conviene ensayarlo con cronómetro: si alguien se alarga, lo primero que se puede recortar está marcado como \emph{opcional}.

\subsection*{Reparto}

\begin{center}
\begin{tabular}{llcl}
\toprule
Integrante & Diapositivas del pptx & Tiempo aprox. & Tema \\
\midrule
Integrante 1 & 1 Portada, 2 Experimento, 3 Implementación & 3:00 & contexto, ventana y linealidad \\
Integrante 2 & 4 Precisión y memoria, 5 DDoS, 6 Scan & 2:30 & resultados de frecuencia \\
Integrante 3 & 7 $\Delta f$ DDoS, 8 $\Delta f$ Scan, 9 Conclusiones & 2:30 & cambio entre ventanas y cierre \\
\midrule
 & \textbf{Total hablado} & \textbf{$\approx$ 8:00} & margen de $\approx$ 2 min \\
\bottomrule
\end{tabular}
\end{center}

El Integrante 1 lleva la diapositiva más importante (la 3), por eso tiene algo más de tiempo. En las preguntas, cada uno responde sobre su parte (sección~\ref{sec:preguntas}).

\begin{center}
\begin{tabular}{clcl}
\toprule
\# & Título en el pptx & Tiempo & Quién \\
\midrule
1 & CMS + CountSketch & 0:30 & Integrante 1 \\
2 & Trazas y criterio de detección & 0:45 & Integrante 1 \\
3 & Ventana deslizante por linealidad & 1:45 & Integrante 1 \\
4 & Precisión y memoria & 0:45 & Integrante 2 \\
5 & DDoS: CMS y CS igualan la detección exacta & 0:45 & Integrante 2 \\
6 & Scan: el ancho pequeño produce falsos positivos & 1:00 & Integrante 2 \\
7 & DDoS: $\Delta f$ marca inicio y salida del ataque & 0:40 & Integrante 3 \\
8 & Scan: CMS-mediana no domina en todos los extremos & 0:40 & Integrante 3 \\
9 & Conclusiones & 1:10 & Integrante 3 \\
\bottomrule
\end{tabular}
\end{center}

\section{Guion por diapositiva}

% -----------------------------------------------------------------------------
\subsection{Integrante 1 — Contexto y diseño (diapositivas 1 a 3, $\approx$ 3:00)}

\subsubsection*{Diapositiva 1 — «CMS + CountSketch» (0:30)}
\textbf{En pantalla:} título, subtítulo «Ventanas deslizantes para detectar ataques en tráfico de red» y nombres.
\decir{Buenas tardes, somos [nombres]. La pregunta de la tarea es cómo saber cuántos paquetes tuvo una IP en los últimos 60 segundos sin guardar todo el tráfico ni recorrerlo de nuevo. Con eso se detecta un ataque cuando una IP concentra demasiado tráfico. Comparamos dos sketches: Count-Min y CountSketch.}

\subsubsection*{Diapositiva 2 — «Trazas y criterio de detección» (0:45)}
\textbf{En pantalla:} traza MAWI, las tarjetas DDoS y Scan, y el cuadro de parámetros.
\decir{Usamos una traza real de MAWI de 2018: 123 millones de paquetes en 15 minutos. Le inyectamos dos ataques de 30 segundos desde el segundo 300. En el DDoS, 4\,000 fuentes atacan una víctima, así que contamos por IP de destino. En el scan, una IP contacta 60\,000 destinos, así que contamos por IP de origen. La ventana es de 60 segundos y avanza cada 10. Una IP es heavy hitter si supera el 1\,\% de la ventana. Probamos cuatro anchos, de 256 a 16\,384.}

\subsubsection*{Diapositiva 3 — «Ventana deslizante por linealidad» (1:45) — \emph{la central}}
\textbf{En pantalla:} el anillo S1…S6, la fórmula $A = S_1 + \dots + S_6$, la regla de rotación y el costo $2\cdot d\cdot w$.
\decir{Este es el diseño central. Dividimos la ventana de 60 segundos en seis subventanas de 10, cada una con su propio sketch en un anillo de seis ranuras. Además mantenemos un agregado A, la suma de las seis. Cada paquete se suma en su ranura y en A, así A siempre representa la ventana completa.
Aquí usamos la linealidad: el sketch de una unión de paquetes es la suma de sus sketches. Entonces, cuando la ventana avanza, restamos de A la ranura más antigua, la limpiamos y la reutilizamos para la nueva. Eso cuesta 2 por d por w contadores, sin importar si la ventana tiene mil u ocho millones de paquetes. Con la misma resta obtenemos gratis el cambio entre ventanas.
Para validar, el total de paquetes de cada ventana se compara con el programa exacto: coincidieron las 84 ventanas. CMS y CS usan exactamente la misma ventana; sólo cambia cómo se actualiza y cómo se estima.}
\textbf{Opcional si sobra tiempo:} los tiempos son enteros en microsegundos y cada subventana es abierta por la izquierda y cerrada por la derecha, igual que \code{exact\_hh}.

\textbf{Traspaso:} \textit{«Ahora [Integrante 2] muestra qué tan precisos son.»}

% -----------------------------------------------------------------------------
\subsection{Integrante 2 — Resultados de frecuencia (diapositivas 4 a 6, $\approx$ 2:30)}

\subsubsection*{Diapositiva 4 — «Precisión y memoria» (0:45)}
\textbf{En pantalla:} error relativo medio de las claves de frecuencia alta sin ataque, y la memoria de 35\,KiB a 2\,240\,KiB.
\decir{Primero validamos sin ataque. Para las IPs más frecuentes, el error relativo de CMS baja de 3,2 con ancho 256 a 0,035 con 16\,384, y el de CS de 0,85 a 0,003. CMS sólo puede sobreestimar, porque todas las colisiones suman. CS usa signos aleatorios: las colisiones se cancelan y el error, que puede ir hacia ambos lados, es menor. Ambos usan la misma memoria, de 35 KiB a 2,2 MiB.}

\subsubsection*{Diapositiva 5 — «DDoS: CMS y CS igualan la detección exacta» (0:45)}
\textbf{En pantalla:} \code{fig\_ddos.png}; abajo, «Latencia: 10\,s … sin falsos positivos».
\decir{En el DDoS, el conteo exacto detecta a la víctima en 310 segundos, 10 después del inicio. CMS y CS la detectan en esa misma evaluación con todos los anchos y sin falsos positivos. Los 10 segundos son la resolución de la ventana. Incluso con ancho 256 el error es bajo: 7,6\,\% en CMS y 3,9\,\% en CS. Al terminar el ataque la curva baja de a poco, a medida que esos paquetes salen de la ventana.}

\subsubsection*{Diapositiva 6 — «Scan: el ancho pequeño produce falsos positivos» (1:00)}
\textbf{En pantalla:} \code{fig\_scan.png}; abajo, «w=256: FP en 310 y 380 s · desde w=1\,024: latencia 20 s y 0 FP».
\decir{El scan es más difícil. El conteo exacto lo detecta en 320 segundos, con latencia de 20. Con ancho 256, ambos sketches lo declaran antes, en 310, y vuelven a dar un falso positivo en 380. En esos instantes la frecuencia real está cerca del umbral y el error de colisión basta para cruzarlo. Desde ancho 1\,024 las decisiones coinciden con la referencia: 20 segundos y cero falsos positivos. Y CS sigue siendo más preciso: con 1\,024, 0,13\,\% de error contra 0,54\,\% de CMS.}

\textbf{Traspaso:} \textit{«[Integrante 3] muestra qué aporta mirar el cambio entre ventanas.»}

% -----------------------------------------------------------------------------
\subsection{Integrante 3 — $\Delta f$ y conclusiones (diapositivas 7 a 9, $\approx$ 2:30)}

\subsubsection*{Diapositiva 7 — «DDoS: $\Delta f$ marca inicio y salida del ataque» (0:40)}
\textbf{En pantalla:} \code{fig\_delta\_ddos.png}; abajo, «$\Delta f$ exacto: +100\,008 en 310\,s · −100\,007 en 370\,s».
\decir{La frecuencia acumula 60 segundos y reacciona lento. El cambio entre ventanas es más nítido: sube 100\,000 paquetes en 310 segundos, cuando entra el ataque, y baja lo mismo en 370, cuando sale. CS lo estima con su estimador normal porque ya maneja signos. En CMS el mínimo no sirve con negativos, así que usamos la mediana, una variante experimental. Con ancho 16\,384 ambos quedan a 2 paquetes o menos del valor exacto.}

\subsubsection*{Diapositiva 8 — «Scan: CMS-mediana no domina en todos los extremos» (0:40)}
\textbf{En pantalla:} \code{fig\_delta\_scan.png}; abajo, «$\Delta f$ exacto: +80\,000 en 310\,s · −80\,000 en 370\,s».
\decir{En el scan el salto es de 80\,000. Con ancho 256 no hay ganador claro: CMS-mediana estimó mejor el incremento, 78 paquetes de error contra 133 de CS, pero peor el decremento, 417 contra 155. Con anchos grandes ambos convergen. La diferencia de fondo es que CS tiene garantías para valores con signo y la CMS-mediana no.}

\subsubsection*{Diapositiva 9 — «Conclusiones» (1:10)}
\textbf{En pantalla:} las cuatro conclusiones numeradas y la pregunta de cierre.
\decir{Cuatro conclusiones. Uno: la linealidad hace viable la ventana deslizante; cada avance cuesta 2dw contadores, sin importar el tráfico. Dos: al aumentar el ancho bajan los errores. Con la misma memoria, CS fue más preciso que CMS en todos los anchos y en ambos ataques; CMS conserva un sesgo positivo por colisiones. Tres: el DDoS se detecta igual que con el conteo exacto con cualquier ancho; el scan con ancho 256 da falsos positivos y necesita al menos 1\,024, unos 140 KiB. Cuatro: Delta f muestra cuándo entra el ataque y cuándo sale de la ventana. En resumen, CS con un ancho de 1\,024 a 4\,096 da las mismas decisiones que el conteo exacto con menos de 1 MB. Gracias, quedamos atentos a sus preguntas.}

% -----------------------------------------------------------------------------
\section{Preguntas probables y respuestas cortas}\label{sec:preguntas}

Entre paréntesis, quién responde por defecto.

\begin{description}[style=nextline]
    \item[¿Por qué no basta con un solo sketch que se reinicia cada 60 segundos? (Integrante 1)]
    Sería una ventana fija, no deslizante: justo después de reiniciar no habría historia. Con el anillo, la ventana siempre contiene los últimos 60\,s.

    \item[¿Cuánta memoria extra cuesta el anillo? (Integrante 1)]
    Siete sketches en vez de uno ($m+1$), más \code{dA}. Es el precio de poder restar lo que expira sin guardar los paquetes.

    \item[¿Cómo saben que sus ventanas coinciden con las del programa exacto? (Integrante 1)]
    El programa compara $N_j$ con \code{exact\_hh} en cada ventana y aborta si difiere; pasó en las 84 ventanas de todas las corridas.

    \item[¿Por qué CMS tiene falsos positivos y nunca falsos negativos? (Integrante 2)]
    Su estimación es siempre mayor o igual que la real: si la real cruza el umbral, la estimada también.

    \item[¿Por qué la latencia es 10\,s en el DDoS y 20\,s en el scan? (Integrante 2)]
    Se evalúa cada 10\,s. El DDoS envía más paquetes por segundo y supera el 1\,\% en la primera evaluación; el scan necesita dos.

    \item[¿Qué pasa con ataques de menor intensidad? (Integrante 2)]
    Lo probamos con 3\,000 paquetes/s (está en el informe, no en el pptx): ambos sketches detectan con latencia de 30\,s, igual que el exacto, sin falsos positivos. El error crece: con $w=256$, 25\,\% en CMS y 13\,\% en CS para el DDoS.

    \item[¿Los resultados dependen de la semilla del hash? (Integrante 2)]
    Repetimos con semillas 7 y 1234. La tendencia se mantiene; el error de CS fluctúa más entre semillas.

    \item[¿Por qué CS se trunca a cero? (Integrante 3)]
    La mediana con signo puede dar negativo y una frecuencia no puede serlo, así que se reporta $\max(0,\hat f)$. Para $\Delta f$ no se trunca.

    \item[¿Por qué la CMS-mediana pierde las garantías del CMS? (Integrante 3)]
    La garantía del CMS depende de que los contadores sean no negativos y de tomar el mínimo. En $\Delta A$ hay valores negativos, así que el mínimo deja de acotar el error y la mediana no tiene esa garantía.
\end{description}

\end{document}
